\section{Testes e Validação}
\label{sec:testes-validacao}

\subsection{Pytest}

A validação do módulo foi realizada com a utilização de Pytest, uma ferramenta que permitiu estruturar testes automatizados para verificar o comportamento das funcionalidades implementadas. Esse processo foi essencial para confirmar que as regras de negócio e os fluxos principais continuavam funcionando corretamente.

\subsection{Pré-configuração}

Antes da execução dos testes, foi necessária a configuração do ambiente, incluindo dependências, dados iniciais e parametrização básica do projeto. Essa etapa garantiu que os testes fossem executados em um cenário consistente com o ambiente real de desenvolvimento.

\subsection{Testes unitários}

Os testes unitários foram empregados para validar as partes mais isoladas da lógica, como validações de regras, criação de pedidos e atualização de registros. Essa abordagem aumentou a confiabilidade da implementação e reduziu a chance de erros em componentes específicos.

\subsection{Testes de integração}

Os testes de integração verificaram a comunicação entre os módulos e a correta interação entre backend, banco de dados e interfaces do sistema. Essa camada de validação foi importante para confirmar que as funcionalidades operavam de forma coerente no contexto completo da aplicação.

\subsection{Processo de Integração Contínua}

O processo de integração contínua foi adotado como mecanismo de validação recorrente, permitindo automatizar a execução dos testes e reduzir erros durante o desenvolvimento. Esse fluxo contribui para manter a qualidade do código e garantir que novas alterações não comprometam o funcionamento do sistema.
